%%% FIELD proof-general-frontier
By \texttt{ass:ordinal-support}, \(L\geq 3\), so \(X\) and \(E\) are finite.  By \texttt{ass:coverage-level} and the definitions of \(c\) and \(\tau\),
\[
0<c<1,\qquad \tau=1+c.
\tag{G1}
\]
Fix \(k\in\{1,\ldots,2L-1\}\).  The rectangle normal form in \texttt{thm:all-minimizers} gives, for every \(S\subseteq E\),
\[
S\in\mathcal V_{\mu,c}
\quad\Longleftrightarrow\quad
\text{there is }(A,B)\in\mathcal F_\tau\text{ with }A-B\subseteq S.
\tag{G2}
\]
Suppose first that a valid \(S\) satisfies \(|S|\leq k\).  Choose \((A,B)\) from the right-hand side of \((\mathrm{G2})\).  Then \(|A-B|\leq |S|\leq k\), while \texttt{def:feasible-rectangles} gives
\[
\mu_1(A)+\mu_0(B)\geq\tau.
\tag{G3}
\]
Conversely, if nonempty \(A,B\subseteq X\) obey \(|A-B|\leq k\) and \((\mathrm{G3})\), then \((A,B)\in\mathcal F_\tau\), and \((\mathrm{G2})\), applied with \(S=A-B\), makes \(A-B\) a valid set of cardinality at most \(k\).  The feasible family in the definition of \(R_k\) is finite and nonempty, since it contains every singleton--singleton pair.  Thus its maximum is attained, and \texttt{def:frontier-profiles} yields
\[
\bigl[\exists S\in\mathcal V_{\mu,c}:|S|\leq k\bigr]
\quad\Longleftrightarrow\quad R_k\geq\tau.
\tag{G4}
\]

We next prove the interval assertion by an integer-hull argument.  Suppose \(B_k\geq\tau\).  All maxima in \texttt{def:concentration-profiles} and \texttt{def:frontier-profiles} range over finite nonempty sets.  Hence there are \(a,b\in\{1,\ldots,L\}\) with \(a+b\leq k+1\), a consecutive \(A_0\subseteq X\) of size \(a\), and a consecutive \(B_0\subseteq X\) of size \(b\), such that
\[
\mu_1(A_0)+\mu_0(B_0)=Q_1(a)+Q_0(b)=B_k\geq\tau.
\tag{G5}
\]
Direct subtraction of two integer intervals gives
\[
A_0-B_0
=\{\min A_0-\max B_0,\ldots,\max A_0-\min B_0\},
\tag{G6}
\]
and therefore
\[
|A_0-B_0|=a+b-1\leq k.
\tag{G7}
\]
By \((\mathrm{G5})\) and \texttt{def:feasible-rectangles}, \((A_0,B_0)\in\mathcal F_\tau\).  Equations \((\mathrm{G2})\)--\((\mathrm{G7})\) show that \(A_0-B_0\) is a valid effect interval of cardinality at most \(k\).

Conversely, suppose that \(I\in\mathcal I_L\) is valid and \(|I|\leq k\).  By \((\mathrm{G2})\), choose \((A,B)\in\mathcal F_\tau\) with \(A-B\subseteq I\).  The sets \(A\) and \(B\) are nonempty by \texttt{def:feasible-rectangles}, so their consecutive hulls
\[
\overline A=\{\min A,\min A+1,\ldots,\max A\},
\qquad
\overline B=\{\min B,\min B+1,\ldots,\max B\}
\tag{G8}
\]
are well defined.  Both endpoint differences \(\min A-\max B\) and \(\max A-\min B\) belong to \(A-B\), hence to \(I\).  Since \(I\) is consecutive, it contains every integer between them.  Thus
\[
\overline A-\overline B
=\{\min A-\max B,\ldots,\max A-\min B\}
\subseteq I.
\tag{G9}
\]
Put \(a=|\overline A|\) and \(b=|\overline B|\).  From \((\mathrm{G9})\),
\[
a+b-1=|\overline A-\overline B|\leq |I|\leq k.
\tag{G10}
\]
All coordinates of \(\mu_1\) and \(\mu_0\) are nonnegative because these vectors lie in \(\Delta_L\).  Since \(A\subseteq\overline A\), \(B\subseteq\overline B\), and \((A,B)\in\mathcal F_\tau\), \texttt{def:concentration-profiles} gives
\[
Q_1(a)+Q_0(b)
\geq \mu_1(\overline A)+\mu_0(\overline B)
\geq \mu_1(A)+\mu_0(B)
\geq\tau.
\tag{G11}
\]
The pair \((a,b)\) is feasible in the maximum defining \(B_k\) by \((\mathrm{G10})\), so \((\mathrm{G11})\) implies \(B_k\geq\tau\).  We have proved
\[
\bigl[\exists I\in\mathcal I_L:h_\mu(I)\geq c,\ |I|\leq k\bigr]
\quad\Longleftrightarrow\quad B_k\geq\tau.
\tag{G12}
\]

The full effect support \(E\) is valid: every \(\gamma\in\Gamma(\mu_1,\mu_0)\) has total mass one by \texttt{def:coupling-polytope}, so \texttt{def:worst-coverage} gives \(h_\mu(E)=1\geq c\) by \((\mathrm{G1})\).  Consequently the minima in \texttt{def:shortest-family} and \texttt{def:interval-optimum} exist.  From \((\mathrm{G4})\) and \((\mathrm{G12})\),
\[
k^*=\min\{k:R_k\geq\tau\},
\qquad
I^*=\min\{k:B_k\geq\tau\}.
\tag{G13}
\]
If \(B_k<\tau\leq R_k\), then \((\mathrm{G13})\) implies \(k^*\leq k<I^*\).  Conversely, if \(k^*<I^*\), then \((\mathrm{G13})\), evaluated at \(k=k^*\), gives \(B_{k^*}<\tau\leq R_{k^*}\).  Hence
\[
k^*<I^*
\quad\Longleftrightarrow\quad
\exists k\in\{1,\ldots,2L-1\}:B_k<\tau\leq R_k.
\tag{G14}
\]

Under nearest-neighbor adjacency on the integer set \(E\), a nonempty subset is connected if and only if it is a consecutive interval.  If \(k^*<I^*\), no member of \(\mathcal K^*\) can be an interval, because such a member would be a valid interval of cardinality \(k^*<I^*\).  Thus every member of \(\mathcal K^*\) is disconnected.  Conversely, if every member of \(\mathcal K^*\) is disconnected but \(k^*=I^*\), an interval attaining the finite minimum \(I^*\) would be a member of \(\mathcal K^*\), a contradiction.  Therefore
\[
k^*<I^*
\quad\Longleftrightarrow\quad
\text{every member of }\mathcal K^*\text{ is disconnected}.
\tag{G15}
\]

It remains to verify both directions of the assertion about the deterministic selector.  The lexicographic order in \texttt{def:deterministic-selector} compares cardinality first, so \(S^\star\in\mathcal K^*\).  Thus the right-hand side of \((\mathrm{G15})\) implies that \(S^\star\) is disconnected.  For the converse, suppose \(S^\star\) is disconnected and some \(J\in\mathcal K^*\) is an interval.  Both sets have \(k^*\) points.  An interval with \(k^*\) integer points has
\[
\operatorname{span}(J)=k^*-1,
\tag{G16}
\]
whereas a disconnected \(k^*\)-point integer set omits at least one integer strictly between its endpoints and hence has
\[
\operatorname{span}(S^\star)+1\geq k^*+1,
\qquad
\operatorname{span}(S^\star)\geq k^*.
\tag{G17}
\]
Equations \((\mathrm{G16})\)--\((\mathrm{G17})\) say that \(J\prec S^\star\): cardinality ties, and \(J\) has strictly smaller span, so the binary code is never consulted.  This contradicts \(S^\star=\min_\prec\mathcal V_{\mu,c}\) in \texttt{def:deterministic-selector}.  Therefore no member of \(\mathcal K^*\) is an interval.  Combining this fact with \((\mathrm{G14})\)--\((\mathrm{G15})\) proves all claimed equivalences, including the selector equivalence at boundary marginal laws.

%%% FIELD proof-ternary-frontier
Assume \(L=3\), as in the statement, and write
\[
p_i=\mu_{1,i},\qquad q_j=\mu_{0,j},\qquad
m_1=\max_{0\leq i\leq2}p_i,qquad m_0=\max_{0\leq j\leq2}q_j.
\tag{T1}
\]
For every \(C\subseteq X\), use the shorthand \(p(C)=\sum_{i\in C}p_i\) and \(q(C)=\sum_{j\in C}q_j\).  The simplex restriction on \(\mu_1,\mu_0\) makes all \(p_i,q_j\) nonnegative and gives \(\sum_i p_i=\sum_jq_j=1\).  We first enumerate all possible difference-set supports; this is the finite-support combinatorial step.

Let \(A=\{a_1<\cdots<a_r\}\) and write \(-B=\{u_1<\cdots<u_s\}\).  The chain
\[
a_1+u_1<\cdots<a_r+u_1<a_r+u_2<\cdots<a_r+u_s
\tag{T2}
\]
consists of \(r+s-1\) distinct members of \(A-B=A+(-B)\).  Consequently,
\[
|A-B|\geq |A|+|B|-1.
\tag{T3}
\]
In particular, if \(|A-B|\leq2\), the only factor-size patterns are \((1,1)\), \((1,2)\), and \((2,1)\).  Each two-point subset of \(X=\{0,1,2\}\) is either adjacent or equal to \(\{0,2\}\).

For an adjacent control factor, its mass is either \(1-q_0\) or \(1-q_2\).  Hence the largest Hall excess \(p(A)+q(B)-1\) among singleton--adjacent rectangles is
\[
m_1-\min(q_0,q_2).
\tag{T4}
\]
Symmetrically, the largest excess among adjacent--singleton rectangles is
\[
m_0-\min(p_0,p_2).
\tag{T5}
\]
The largest singleton--singleton excess is \(m_1+m_0-1\).  It is dominated by \((\mathrm{T4})\) because \(\min(q_0,q_2)\leq1-m_0\), and it is also dominated by \((\mathrm{T5})\) because \(\min(p_0,p_2)\leq1-m_1\).  Applying the rectangle formula \texttt{lem:hall-coverage} and the definition \texttt{def:ternary-thresholds}, we therefore obtain
\[
\max_{\substack{I\in\mathcal I_3\\ |I|\leq2}}h_\mu(I)
=\max\{m_1-\min(q_0,q_2),\ m_0-\min(p_0,p_2),\ 0\}
=b_2.
\tag{T6}
\]

The gapped two-point factor rectangles are precisely
\[
\{i\}\times\{0,2\}
\quad\text{and}\quad
\{0,2\}\times\{j\},
\qquad i,j\in X.
\tag{T7}
\]
Their respective difference sets and Hall excesses are
\[
\{i\}-\{0,2\}=\{i-2,i\},
\qquad p_i+q_0+q_2-1=p_i-q_1,
\tag{T8}
\]
and
\[
\{0,2\}-\{j\}=\{-j,2-j\},
\qquad p_0+p_2+q_j-1=q_j-p_1.
\tag{T9}
\]
Thus the largest excess furnished by a gapped factor is
\[
\max\{m_1-q_1,\ m_0-p_1,\ 0\}=d_2.
\tag{T10}
\]
Equations \((\mathrm{T3})\)--\((\mathrm{T10})\) exhaust every rectangle whose difference set is contained in a two-point effect set.  Notice that a singleton rectangle may also lie in a disconnected pair; when \(c>b_2\), equation \((\mathrm{T6})\) makes every such singleton certificate strictly smaller than \(c\), so only \((\mathrm{T7})\) can certify a disconnected pair at level \(c\).

We next classify difference sets of cardinality three or more.  By \((\mathrm{T3})\), a new three-point support can only have factor sizes \((1,3)\), \((3,1)\), or \((2,2)\).  The first two supports are consecutive intervals.  In the \((2,2)\) case, write the positive gaps of the two factors as \(g_A,g_B\in\{1,2\}\).  Of the four formal differences, the two interior ones coincide exactly when \(g_A=g_B\).  If \(g_A=g_B=1\), the resulting three points are consecutive.  If \(g_A=g_B=2\), both factors must be \(\{0,2\}\), and the resulting support is the sole disconnected three-point difference set
\[
T=\{0,2\}-\{0,2\}=\{-2,0,2\}.
\tag{T11}
\]
If \(g_A\neq g_B\), the four differences are distinct and, since \(\{g_A,g_B\}=\{1,2\}\), they are four consecutive integers.  If one factor is all of \(X\) and the other has two points, an adjacent two-point factor gives four consecutive differences and the factor \(\{0,2\}\) gives all five differences.  If both factors equal \(X\), the difference set is all of \(E\).  Together with the singleton-factor cases, this proves that every four- or five-point difference set is an interval and that \(T\) is the only disconnected three-point difference set.  The enumeration used only \(X=\{0,1,2\}\), not positivity of any cell.

For \(d\in\{0,1\}\), \texttt{def:concentration-profiles} gives
\[
Q_d(1)=m_d,qquad
Q_d(2)=1-\min(\mu_{d,0},\mu_{d,2}),
\qquad Q_d(3)=1.
\tag{T12}
\]
For \(k=3\), the undominated size pairs in \texttt{def:frontier-profiles} are \((1,3)\), \((3,1)\), and \((2,2)\); the remaining pairs have no larger profile value because each \(Q_d\) is nondecreasing in its width.  Hence \texttt{def:ternary-thresholds} and \((\mathrm{T12})\) give
\[
B_3-1
=\max\{m_1,m_0,1-\min(p_0,p_2)-\min(q_0,q_2)\}
=b_3.
\tag{T13}
\]
By the interval equivalence in \texttt{thm:general-frontier} and \(\tau=1+c\), equation \((\mathrm{T13})\) says
\[
\text{an interval of cardinality at most three is valid}
\quad\Longleftrightarrow\quad b_3\geq c.
\tag{T14}
\]
Also, directly from \texttt{def:ternary-thresholds},
\[
d_2\leq\max(m_1,m_0)\leq b_3.
\tag{T15}
\]

It remains to evaluate the exceptional support \(T\).  A rectangle \(A\times B\) has \(A-B\subseteq T\) exactly when all its row and column indices have the same parity.  Since \(X=\{0,1,2\}\), every nonempty such rectangle lies either in \(\{0,2\}\times\{0,2\}\) or in \(\{1\}\times\{1\}\).  Nonnegativity of the marginal cells means that the largest excess in the even block is attained by the whole even block, while the odd block has only one nonempty rectangle.  Therefore \texttt{lem:hall-coverage} yields
\[
h_\mu(T)
=\max\{0,\ (p_0+p_2)+(q_0+q_2)-1,\ p_1+q_1-1\}
=\max\{0,\ 1-p_1-q_1,\ p_1+q_1-1\}.
\tag{T16}
\]
If \(c>b_3\), then \((\mathrm{T13})\) gives \(p_1\leq m_1<c\) and \(q_1\leq m_0<c\).  Since \(c<1\) by \texttt{ass:coverage-level},
\[
p_1+q_1-1<2c-1<c.
\tag{T17}
\]
Thus, in the regime \(c>b_3\), equations \((\mathrm{T16})\)--\((\mathrm{T17})\) show
\[
h_\mu(T)\geq c
\quad\Longleftrightarrow\quad 1-p_1-q_1\geq c.
\tag{T18}
\]

We now prove sufficiency of the two displayed regimes.  Suppose first that
\[
d_2\geq c>b_2.
\tag{T19}
\]
By \((\mathrm{T8})\)--\((\mathrm{T10})\), a gapped two-point difference set has a rectangle excess at least \(c\), so it is valid by \texttt{lem:hall-coverage}.  Equation \((\mathrm{T6})\) and \(c>b_2\) exclude every valid interval of cardinality at most two, including every singleton.  Therefore \(k^*=2\), and every member of \(\mathcal K^*\) is disconnected.  The exact all-minimizer identity in \texttt{thm:all-minimizers}, together with the exhaustive patterns \((\mathrm{T7})\)--\((\mathrm{T10})\), gives
\[
\mathcal K^*
=\bigl\{\{i-2,i\}:p_i+q_0+q_2-1\geq c,\ i\in X\bigr\}
\mathbin{\cup}
\bigl\{\{-j,2-j\}:q_j+p_0+p_2-1\geq c,\ j\in X\bigr\},
\tag{T20}
\]
where the union is a family of sets and hence automatically deduplicates repeated supports.

Suppose instead that
\[
1-p_1-q_1\geq c>b_3.
\tag{T21}
\]
Equations \((\mathrm{T16})\)--\((\mathrm{T18})\) make \(T\) valid.  Equation \((\mathrm{T14})\) excludes every valid interval of size at most three, and \((\mathrm{T15})\) gives \(d_2<c\), excluding the only remaining two-point certificates after \((\mathrm{T6})\).  Hence no set of cardinality at most two is valid, \(k^*=3\), and the support enumeration forces
\[
\mathcal K^*=\{T\}=\bigl\{\{-2,0,2\}\bigr\}.
\tag{T22}
\]

For necessity, assume every member of \(\mathcal K^*\) is disconnected.  The full set \(E\) is valid, as shown from \texttt{def:coupling-polytope} and \texttt{def:worst-coverage} in the proof of \texttt{thm:general-frontier}, so \(1\leq k^*\leq5\).  A one-point set is connected.  Moreover, \texttt{thm:all-minimizers} says that every member of \(\mathcal K^*\) is a difference set, and the support enumeration above says that every four- or five-point difference set is connected.  Thus \(k^*\in\{2,3\}\).

If \(k^*=2\), some feasible rectangle generates a disconnected two-point member of \(\mathcal K^*\).  The classification \((\mathrm{T7})\)--\((\mathrm{T10})\) then gives \(d_2\geq c\).  No interval of cardinality at most two can be valid, since it would also belong to \(\mathcal K^*\); by \((\mathrm{T6})\), this is equivalent to \(c>b_2\).  Hence \((\mathrm{T19})\) holds.

If \(k^*=3\), the only possible disconnected difference set is \(T\), by \((\mathrm{T11})\), so \(\mathcal K^*=\{T\}\).  No interval of cardinality at most three is valid, since it would be another member of \(\mathcal K^*\); equation \((\mathrm{T14})\) therefore gives \(c>b_3\).  The validity of \(T\), followed by \((\mathrm{T18})\), gives \(1-p_1-q_1\geq c\).  Thus \((\mathrm{T21})\) holds.

The two regimes are mutually exclusive: the first requires \(d_2\geq c\), whereas the second and \((\mathrm{T15})\) imply \(d_2\leq b_3<c\).  Finally, \texttt{thm:general-frontier} proves that every member of \(\mathcal K^*\) is disconnected if and only if the span-tie-broken selector \(S^\star\) is disconnected.  All feasibility tests above use weak inequalities \(\geq\), and every exclusion uses the logically complementary strict inequality \(>\).  Since the rectangle normal form in \texttt{thm:all-minimizers} permits zero atoms and does not require uniqueness, the classification includes zero cells, exact threshold equalities, repeated difference sets, and all span and code ties.
